% !TITLE 终南别业
% !AUTHOR 王维
% !DYNASTY 唐
% !GENRE 五言律诗
% !ORDER 39

\begin{poem}
中岁颇好道\textsuperscript{1}，晚家南山陲\textsuperscript{2}。
兴来每独往，胜事空自知\textsuperscript{3}。
行到水穷处\textsuperscript{4}，坐看云起时\textsuperscript{5}。
偶然值林叟\textsuperscript{6}，谈笑无还期\textsuperscript{7}。
\end{poem}

\begin{annotations}
\notetext{1}{中岁颇好道：中年以后很喜欢佛道。好（hào），喜好；道，佛理禅机。}
\notetext{2}{南山陲：终南山的边陲地带。陲（chuí），边。晚年把家安在了终南山脚下。}
\notetext{3}{胜事空自知：美好的感受只有自己知道。无人分享，也不想分享——自足自乐。}
\notetext{4}{行到水穷处：沿着溪流走到水的尽头。穷，尽。走投无路时，坐下来——不焦虑，不回头。}
\notetext{5}{坐看云起时：坐下来看云升起。水的尽头是云的开始，穷和起之间是一种自然的转换。}
\notetext{6}{值林叟：遇到山林中的老人。值，偶遇；叟（sǒu），老翁。}
\notetext{7}{无还期：没有回去的时限。谈得高兴了就忘了时间——隐居生活中最美的就是这种没有计划的邂逅。}
\end{annotations}

\begin{appreciation}
这首诗写于王维晚年隐居终南山辋川别业之时。中岁以后的王维，经历了安史之乱的屈辱，看透了官场的险恶，最终选择了归隐。这首诗记录的，便是他隐居生活中的一个寻常日子。

中岁颇好道，晚家南山陲——中年以后颇喜好佛道之道，晚年把家安在了终南山的边陲。陲（chuí）就是边。好道，是王维人生的转折点。他从一个意气风发的少年才子，变成了一个淡泊宁静的山中隐士。这个转变不是突然的，是中岁以后一点一点发生的。

兴来每独往，胜事空自知——兴致来了就独自出游，美好的感受只有自己知道。独往，是独自前往，没有同伴。空自知，是只有自己知道，无法与人分享。这不是孤独，是自足。有些美好，确实无法分享，也不需要分享。

行到水穷处，坐看云起时——走到水的尽头，便坐下来看云升起。这是全诗最著名的一句，也是王维人生哲学的最好写照。水穷处，是路的尽头，也是困境的象征。但王维不焦虑，不回头，他只是坐下来，看云起。云从山间升起，变幻莫测，无始无终。水的尽头，正是云的开始。穷和起之间，是一种自然的过渡，也是一种人生的顿悟。

偶然值林叟，谈笑无还期——偶然遇到山林中的老人，谈笑风生，忘记了回去的时间。叟（sǒu）就是老人。值是遇到。无还期，是没有回去的期限，也就是不想回去了。这种偶遇，这种忘归，是隐居生活最迷人的地方：没有计划，没有目的，一切都是偶然，一切又都恰到好处。

行到水穷处，坐看云起时。这十个字，道尽了中国文人最理想的生活态度：不执着，不焦虑，随遇而安，顺其自然。水尽了，有云；云散了，有风。世间万物，循环往复，何必强求？
\end{appreciation}
